% ATTENTION : modifs faites directement dans moderncv.cls et moderncvheadi.sty pour :
% - ajouter la possibilité d'avoir deux adresses e-mail. 
% - ajouter un espacement entre l'adresse et les autres informations personnelles dans le header.
% - réduire l'espacement avant le nom et le titre dans le header.
\documentclass[11pt,a4paper,sans]{moderncv}
\usepackage[utf8]{inputenc}
\usepackage[T1]{fontenc}
\usepackage{lmodern} % optionnel mais fiable pour le gras
\usepackage[scale=0.85]{geometry}
\usepackage{helvet}
\usepackage[french]{babel}
\usepackage{multicol}
\usepackage{blindtext}
\usepackage[unicode]{hyperref}
\usepackage{amsmath}
\usepackage{enumitem}

% config 
\moderncvstyle{classic}
\moderncvcolor{blue}
\setlength{\hintscolumnwidth}{2.7cm}

\definecolor{link}{rgb}{0.16, 0.32, 0.75}

\newcommand\Colorhref[3][link]{\href{#2}{\small\color{#1}#3}}

% \makeatletter
% \newcommand*{\emailtwo}[1]{\def\@emailtwo{#1}}
% \patchcmd{\makecvtitle}{\emailsymbol\emaillink{\@email}}{%
%   \emailsymbol\emaillink{\@email}%
%   \ifdefined\@emailtwo\\\emailsymbol\emaillink{\@emailtwo}\fi%
% }{}{}
% \makeatother

\nopagenumbers

% Configuration des listes itemize
\setlist[itemize]{leftmargin=1em, itemsep=0.3em, topsep=0.3em, label=-}

% personal info 
\name{Frédérique}{LECOURTIER}
\title{Doctorante\\ \tiny \; \\ \tiny \;}
\address{2 Rue Marie Hamm}{67\,000 STRASBOURG}{FRANCE}
\phone[mobile]{(private)}
\email{frederique.lecourtier@inria.fr}
\emailtwo{frederique.lecourtier@gmail.com}
\social[linkedin]{flecourtier}
\social[github]{flecourtier}
\extrainfo{\faGlobe~\href{https://flecourtier.github.io/}{flecourtier.github.io}}
\photo[110pt][0pt]{Photo_CV.jpg}

%programmation langage
\colorlet{body}{black!80!white}
\newcommand{\cvtag}[1]{%
    \tikz[baseline]\node[anchor=base,draw=body!30,rounded corners,inner xsep=1ex,inner ysep =0.75ex,text height=1.5ex,text depth=.25ex]{#1};
}

% Commande pour lister plusieurs tags avec espacement automatique
\usepackage{xparse}
\NewDocumentCommand{\cvtags}{m}{%
    \foreach \tag in {#1} {\cvtag{\tag}}\\[2mm]%
}

% Commande pour lister plusieurs items avec des points de séparation
\newcommand{\cvsep}{\space$\boldsymbol{\cdot}$\space}
\NewDocumentCommand{\cvitems}{m}{%
    \foreach \item [count=\i, remember=\i as \lastn] in {#1} {%
        \ifnum\i=1 \item \else \cvsep \item \fi%
    }%
}

\begin{document}
    \makecvtitle 
    \vspace*{-2.2cm}
    {\small\emph{(Derni\`ere modification le \today)}}

    \section{Thèmes de recherche}

        \cvitems{Calcul scientifique, Apprentissage automatique guidé par la physique, Analyse numérique, Méthodes des éléments finis, Solveurs hybrides, Équations aux dérivées partielles, Estimation d'erreur}

    \section{Formations}

        \cventry{2023 - ...}{\textbf{Doctorat}\normalfont\ en Mathématiques Appliquées}{Université de Strasbourg}{}{France}{
            \textit{Équipe:} Mimesis, Centre Inria de l'Université de Lorraine. \\
            \textit{Supervisé par:} Emmanuel Franck, Michel Duprez et Vanessa Lleras.
        }
        \cventry{2021 - 2023}{\textbf{Master}\normalfont\ Calcul Scientifique et Mathématiques de l'Information (CSMI)}{Université de Strasbourg}{}{France - Mention Bien}{}
        \cventry{2018 - 2021}{\textbf{Licence}\normalfont\ en Mathématiques Appliquées}{Université de Strasbourg}{}{France - Mention Bien}{}
        \cventry{2018}{\textbf{Baccalauréat}\normalfont\ Scientifique}{Lycée LeCorbusier}{Illkirch}{France}{}{}

    \section{Expériences}

        \cventry{2023 - ...}{Doctorat\normalfont\ en Mathématiques Appliquées}{Mimesis}{Strasbourg}{France}
        {%
            \cvtags{Python, PyTorch, Jax, \Colorhref{https://fenicsproject.org/}{FEniCSx}, \Colorhref{https://www.scimba.org/}{ScimBa}}
            \textit{Sujet:} Méthodes éléments finis et réseaux de neurones pour la chirurgie augmentée. \\
            \textit{Objectifs:}
            \begin{itemize}
                \item Développer des méthodes hybrides combinant les méthodes éléments finis et les réseaux de neurones pour la résolution de problèmes d'équations aux dérivées partielles paramétriques.
                \item Développer des méthodes scientifiques d'apprentissage automatique dans ScimBa (bibliothèque Python en cours de développement dans l'équipe-projet Macaron, Inria) pour la résolution de problèmes d'équations aux dérivées partielles paramétriques.
            \end{itemize}
            \textcolor{link}{\faLink} \Colorhref{https://github.com/flecourtier/EnrichedFEMUsingPINNs}{Dépôt GitHub.}
        }

        \cventry{2023 \\ (6 mois)}{Stage}{Mimesis}{Strasbourg}{France}
        {%
            \cvtags{Python, Tensorflow, PyTorch, \Colorhref{https://olddocs.fenicsproject.org/dolfin/latest/python/}{FEniCS}, GitHub, VSCode}
            \textit{Sujet:} Méthodes éléments finis et réseaux de neurones pour la chirurgie augmentée. \\
            \textit{Supervisé par:} Emmanuel Franck, Michel Duprez et Vanessa Lleras. \\
            % \textit{Objectifs:} Introduction au sujet de thèse. \\
            \textcolor{link}{\faLink} \Colorhref{https://github.com/flecourtier/phifem_stage}{Dépôt GitHub.} \qquad \textcolor{link}{\faLink} \Colorhref{https://csmi.cemosis.fr/csmi-stages-2023/m2/_attachments/Lecourtier-Frédérique.pdf}{Rapport.}
        }
        \cventry{2022 \\ (2 mois)}{Stage}{Cemosis}{Strasbourg}{France}
        {%
            \cvtags{C++, Python, MPI, \Colorhref{https://docs.feelpp.org/home/index.html}{Feel++}, Antora, Sphinx, Doxygen, VSCode, CMake}
            \textit{Sujet:} Méthode Parareal pour la résolution de l'équation de la chaleur en Feel++. \\
            \textit{Supervisé par:} Christophe Prudhomme et Luca Berti. \\
            \textit{Objectifs:}
            \begin{itemize}
                \item Comprendre et implémenter en Python et en C++, la méthode Parareal, méthode d'intégration parallèle en temps.
                \item Implémenter la résolution de l'équation de la chaleur avec Feel++ (bibliothèque C++ développée dans l'équipe Cemosis, IRMA).
            \end{itemize}
            \textcolor{link}{\faLink} \Colorhref{https://github.com/master-csmi/2022-m1-lorenz}{Dépôt GitHub.} \qquad \textcolor{link}{\faLink} \Colorhref{https://csmi.cemosis.fr/csmi-stages-2022/m1/_attachments/Lecourtier-Frédérique.pdf}{Rapport.}
        }

    \section{Publications}

        \cventry{2026}{\normalfont H. Barucq, M. Duprez, F. Faucher, E. Franck, \underline{F. Lecourtier}, V. Lleras, V. Michel-Dansac, N. Victorion. \textbf{Enriching continuous Lagrange finite element approximation spaces using neural networks}. Mathematical Modelling and Numerical Analysis (M2AN).}{}{}{}{
            (\Colorhref{https://doi.org/10.1051/m2an/2026007}{pdf}), (\Colorhref{https://arxiv.org/abs/2502.04947}{pré-publication [arXiv}, \Colorhref{https://hal.science/hal-04935072v4/document}{HAL]})
        }

    \section{Conférences, présentations et posters}

        \cventry{2026\\29 Juin-3 Juillet}{SciCADE 2026}{George Square campus}{Édimbourg}{Royaume-Uni}{
            \textcolor{link}{\faLink} \Colorhref{https://scicade.org/}{Site web.} %\qquad \textcolor{link}{\faLink} \Colorhref{https://flecourtier.github.io/images/talks/2025_07_15.pdf}{Présentation.}
        }

        \cventry{2025\\17 Décembre}{Séminaire doctorant}{Institut IMAG}{Montpellier}{France}{
            \textcolor{link}{\faLink} \Colorhref{https://imag.umontpellier.fr/?page_id=625&idsem=596}{Site web.} \qquad \textcolor{link}{\faLink} \Colorhref{https://flecourtier.github.io/images/talks/2025_12_17.pdf}{Présentation.}
        }

        \cventry{2025\\13-18 Juillet}{ICOSAHOM 2025}{Université McGill}{Montréal}{Canada}{
            \textcolor{link}{\faLink} \Colorhref{https://icosahom2025.org/}{Site web.} \qquad \textcolor{link}{\faLink} \Colorhref{https://flecourtier.github.io/images/talks/2025_07_15.pdf}{Présentation.}
        }

        \cventry{2025\\24-26 Mars}{EMS-TAG-SciML-25}{École polytechnique de Milan}{Milan}{Italie}{
            \textcolor{link}{\faLink} \Colorhref{https://www.mate.polimi.it/events/EMS-TAG-SciML-25/index.php}{Site web.} \qquad \textcolor{link}{\faLink} \Colorhref{https://flecourtier.github.io/images/posters/2025_03_24.pdf}{Poster.}
        }

        \cventry{2025\\17-21 Février}{DTE \& AICOMAS 2025}{Arts et Métiers - ENSAM}{Paris}{France}{
            \textcolor{link}{\faLink} \Colorhref{https://dte_aicomas_2025.iacm.info/organizers}{Site web.} \qquad \textcolor{link}{\faLink} \Colorhref{https://flecourtier.github.io/images/talks/2025_02_20.pdf}{Présentation.}
        }

        \cventry{2024\\28-30 Octobre}{CJC-MA 2024}{École Normale Supérieure de Lyon}{Lyon}{France}{
            \textcolor{link}{\faLink} \Colorhref{https://cjc-ma2024.sciencesconf.org/}{Site web.} \qquad \textcolor{link}{\faLink} \Colorhref{https://flecourtier.github.io/images/posters/2024_10_24.pdf}{Poster.}
        }

        \cventry{2024\\8-12 Juillet}{SciML 2024}{UFR de mathématique et d'informatique}{Strasbourg}{France}{
            \textcolor{link}{\faLink} \Colorhref{https://irma.math.unistra.fr/~micheldansac/SciML2024/index.html}{Site web.} \qquad \textcolor{link}{\faLink} \Colorhref{https://flecourtier.github.io/images/posters/2024_07_08.pdf}{Poster.}
        }

    \section{Expériences d'enseignement}

        \cventry{2025 - 2026}{Mission de vacation}{UFR de mathématique et d'informatique}{Strasbourg}{France}{
            50 h de travaux pratiques en informatique - L1 (Python) / L2 (Python). \\
            14 h de cours intégrés en informatique - L2 (Python).
        }

        \cventry{2025}{Rendez-vous des Jeunes Mathématiciennes et Informaticiennes (RJMI)}{UFR de mathématique et d'informatique}{Strasbourg}{France}{
            Encadrement d'ateliers de recherche pour un groupe de lycéennes dans le cadre des RJMI (événement sur 3 jours).
        }

        \cventry{2024 - 2025}{Mission de vacation}{UFR de mathématique et d'informatique}{Strasbourg}{France}{
            64 h de travaux pratiques en informatique - L1 (Python) / L2 (Python) / L3 (C++).
        }

        \cventry{2023 - 2024}{Mission de vacation}{UFR de mathématique et d'informatique}{Strasbourg}{France}{
            50 h de travaux pratiques en informatique - L2 (Python) / L3 (C++).
        }
        
    \section{Compétences}

        \cvitems{GitHub, Python, C++, C, Rust, LaTeX, PyTorch, JAX, ParaView, Gmsh, FEniCS}

    \section{Référents}

        \textbf{Directeur de thèse:}

        Emmanuel Franck : \emaillink{emmanuel.franck@inria.fr}
        
        \smallskip
        \textbf{Co-encadrants de thèse:}

        Michel Duprez : \emaillink{michel.duprez@inria.fr}

        Vanessa Lleras : \emaillink{vanessa.lleras@umontpellier.fr}

        \smallskip
        \textbf{Collaborateur:}

        Victor Michel-Dansac: \emaillink{victor.michel-dansac@inria.fr}

\end{document}